\title{Byte Latent Transformer: Patches Scale Better Than Tokens}

\begin{abstract}
We present the Byte Latent Transformer ({\textsc BLT}), a novel large language model architecture operating at the byte level, which, for the first time, achieves parity with conventional tokenization-based LLMs in terms of performance at scale while offering notable gains in inference robustness and efficiency. {\textsc BLT} works by aggregating bytes into adaptively determined patches, utilizing these as the core computational units. The segmentation of these patches is guided by measuring the entropy of subsequent bytes, thereby directing increased computational resources and representational capacity to segments characterized by higher data complexity. This work introduces the initial scaling analysis of byte-level language models under constrained FLOP budgets, reaching up to 8B parameters and trained over 4T bytes of raw text. Our findings establish that large-scale training of models directly on byte streams, without predefined vocabularies, is viable. Dynamic selection of lengthy patches where the input sequence is more predictable leads to increased efficiency during both the training and inference phases, coupled with observable gains in both reasoning tasks and long-tailed generalization. Overall, at a constant inference cost, {\textsc BLT} demonstrates substantially superior scaling capabilities compared to token-based approaches, due to the concurrent upscaling of both patch sizes and model capacity.
\end{abstract}

\section{Introduction}
\label{section:intro}

We present the Byte Latent Transformer~(\textbf{BLT}), a novel model architecture that dispenses with the need for tokenization by directly ingesting raw byte sequences. This design achieves parity with traditional tokenization-based systems at scale, while also offering notable gains in efficiency and resilience (\S\ref{sec:robustness}). 
Conventional large language models (llms) are predominantly end-to-end, yet they rely on tokenization—a heuristic method that segments byte sequences into a predetermined token vocabulary—as a pre-training step. This reliance introduces representational biases and results in several limitations, including domain and modality sensitivity~\citep{dagangetting}, vulnerability to noisy inputs~(\S\ref{sec:robustness}), poor capture of orthographic details~\citep{edman2024cute}, and disadvantages for some languages in multilingual contexts~\citep{liang2023xlm,petrov2024language,limisiewicz2024myte}.

Historically, tokenization has played a vital role because training llms directly on raw bytes becomes prohibitively expensive at scale, mainly due to increased sequence lengths~\citep{xue2022byt5}. Earlier research addresses this by introducing more computationally efficient self-attention mechanisms~\citep{el2020characterbert, clark2022canine} or moving toward architectures that eliminate attention altogether~\citep{wang2024mambabyte}~(\S\ref{section:related_work}). Nonetheless, these solutions predominantly benefit the training of \textit{small models}. For large-scale models, the major bottleneck in computational resources lies in the substantial feed-forward network layers executed on each byte, rather than in the attention components.

We introduce a dynamic and learnable approach for forming \textit{patches} of bytes~(\S\ref{section:patching}) in order to optimize computational resource distribution, alongside a novel architecture that integrates both byte-level and patch-level data. In contrast to traditional tokenization, {\textsc BLT} does not impose a predefined vocabulary on patches. Instead, it uses compact, trainable encoder and decoder components to transform arbitrary byte sequences into latent patch representations. Our experiments demonstrate that this method leads to a \textit{greater} efficiency in compute allocation compared to models based on tokenization.

Tokenization-based llms distribute an equal amount of computational resources to every token. This approach sacrifices computational efficiency for model performance, as the heuristics used to compress tokens do not consistently correspond to the prediction complexity of individual tokens. Our architecture is built on the principle that computation should be allocated adaptively, focusing resources where greater prediction difficulty exists. For instance, predicting the ending letters of most words is generally straightforward—these are low-entropy choices—whereas selecting the initial word in a sentence is a more complex and uncertain task that could benefit from a larger transformer. {\textsc BLT} embodies this principle in its architectural design~(\S\ref{section:architecture}), utilizing three transformer modules: a pair of compact byte-level \textit{local models} and a singular, more substantial \textit{latent transformer} for global context~(\autoref{fig:arch_high_level}).

To facilitate dynamic compute allocation and the grouping of bytes into patches, {\textsc BLT} leverages the entropy associated with next-byte prediction, generating context-sensitive patches characterized by a relatively consistent distribution of information.

We introduce the inaugural flop-constrained scaling analysis of byte-level architectures, extending up to 8B parameters and processing 4T bytes of training data, and demonstrate the viability of large-scale, end-to-end training from raw bytes—without the necessity of a predefined tokenization scheme. Notably, {\textsc BLT} achieves parity with Llama 3 in terms of training flop efficiency\footnote{Model computational cost is assessed based on the total number of Floating Point OPerations (flops) required.} while realizing as much as a 50\% reduction in inference flops~(\S\ref{section:scaling}).

Our findings further indicate that leveraging raw bytes as inputs leads to marked gains in modeling capability over the rare and idiosyncratic segments of the data distribution. {\textsc BLT} models exhibit greater resilience to input noise compared to models reliant on tokenization, and they display advanced proficiency in tasks requiring fine-grained character-level comprehension, such as orthographic processing, phonological phenomena, and machine translation for low-resource languages~(\S\ref{sec:robustness}).

Moreover, with {\textsc BLT} models, it is possible to jointly scale up both the model’s parameter count and the patch size while adhering to an unchanged inference flop budget. Employing larger patch sizes yields computational savings on average, which can then be reallocated to expand the global latent transformer component, as its invocation frequency decreases. Through controlled experiments focused on fixed inference flops~(\autoref{fig:fixed_inference_scaling}), we observe significantly superior scaling behavior compared to traditional tokenization-based model designs.

In conclusion, this work presents several key advances: 1) We propose {\textsc BLT}, a byte latent llm framework that adaptively apportions computation to boost flop efficiency, 2) We demonstrate parity in training flops with Llama 3 up to 8B parameters, enabling users to exchange small decreases in evaluation performance for flop efficiency improvements as high as 50\%, 3) The {\textsc BLT} architecture introduces a novel direction in llm scaling, facilitating increases in model size without exceeding a set inference computation cost, and 4) We provide empirical evidence that {\textsc BLT} models exhibit greater resilience to noisy inputs and recognize sub-word patterns often overlooked by token-centric llms. The codebase for training and inference with {\textsc BLT} is publicly available at~\url{https://github.com/facebookresearch/blt}.

\section{Patching: From Individual Bytes to Groups of Bytes}
\label{section:patching}

Dividing a byte sequence into \textit{patches} enables {\textsc BLT} to flexibly adjust computational resources according to the local context. Figure~\ref{fig:patching_types} illustrates a range of strategies for partitioning bytes into patches. In formal terms, a patching function $f_p$ takes a sequence of bytes $\pmb{x}=\{x_i,|i=1,\ldots n\}$ of length $n$ and produces a sequence of $m < n$ patches $\pmb{p}=\{p_j|j=1,\ldots,m\}$, assigning each $x_i$ a value from the set \{0,1\}, where a value of 1 signals the initiation of a new patch. Regardless of whether a model processes data as discrete tokens or patches, the overall computational load hinges on the number of iterations performed by the principal Transformer module. Specifically, within {\textsc BLT}, this equates to the count of patches generated via the chosen patching function for data encoding. As a result, the typical size of a patch—referred to as the \textit{patch size}—serves as the principal determinant of processing expense throughout training and evaluation phases for any given patching approach~(\S\ref{section:flops}).

Subsequently, we present three distinct patching functions: fixed-length patching~(\S\ref{section:static-patch}), whitespace-based patching~(\S\ref{section:white-patch}), and entropy-guided dynamic patching using a compact byte-level language model~(\S\ref{section:dyn-patch}). We conclude by elaborating on incremental patching and clarifying distinctions between the concepts of tokenization and patching~(\S\ref{section:bpe-patch}).

\subsection{Strided Patching Every K Bytes}
\label{section:static-patch}

One of the simplest methods for organizing bytes involves dividing them into patches of fixed length $k$, as utilized in MegaByte~\citep{yu2023megabyte}. This approach, which applies a constant stride, is straightforward to realize both during training and inference phases, and offers a simple way to adjust the typical patch size, thus allowing for direct management of computational resources. Despite these advantages, the fixed-size patching strategy has notable drawbacks. Primarily, it does not adapt computation to match the needs of the data: for example, resources might be squandered on trivial content, such as whitespace in programming code, or conversely, inadequate computational effort may be devoted to complex and information-rich byte segments, like those representing mathematical expressions. Additionally, this strategy can result in inconsistent and context-insensitive splitting of comparable byte sequences, causing, for instance, identical words to be partitioned in varying manners.

\subsection{Space Patching}
\label{section:white-patch}

\citet{slagle2024spacebyte} introduces a straightforward yet effective refinement to strided patching, wherein new patches are initiated after encountering any space-like bytes\footnote{
    Space-like bytes are considered any byte that does not correspond to a Latin letter, numeric digit, or utf-8 continuation byte. Furthermore, each patch must include at least one byte that is not space-like. }—locations which commonly mark the boundaries between linguistic elements in numerous languages. Within Space patching, a latent transformer operation (entailing increased computational cost) is devoted to modeling every word individually. This approach standardizes the patching of words throughout sequences and concentrates computational effort on challenging prediction points, which are frequently found just after spaces. For instance, when given the query ``Who composed the Magic Flute?\uline{\hspace{.6em}}'', predicting the initial byte of the correct answer demands much more from the model than generating subsequent bytes after the letter ``M'', since the initial character substantially limits the set of plausible continuations; thus, producing the answer ``Mozart'' becomes increasingly straightforward with each subsequent character. Nevertheless, space patching presents limitations: it is not universally compatible across all languages and domains, and critically, lacks flexibility in patch size adjustment. To address this, we now present a novel patching technique that builds upon the insight that the initial bytes of words are often the most challenging to predict, while simultaneously offering an intuitive means to adjust patch size.

\subsection{Entropy Patching: Using Next-Byte Entropies from a Small Byte LM}
\label{section:dyn-patch}

Instead of utilizing a heuristic based on rules, such as segmentation by whitespace, we adopt a data-driven methodology to locate next-byte predictions exhibiting high uncertainty. We propose a method called \textit{entropy patching}, which leverages entropy estimates to determine patch delimiters.

A compact byte-level auto-regressive language model is trained on the {\textsc BLT} training corpus, allowing us to calculate the next-byte entropy values according to the language model's output distribution $p_{e}$ across the byte vocabulary $\mathcal{V}$:

\begin{equation}
H(x_i) = \sum_{v \in \mathcal{V}} p_{e}(x_i=v|\pmb{x}_{<i}) \log p_{e}(x_i=v|\pmb{x}_{<i})
\end{equation}

We investigate two approaches for determining patch boundaries using entropies $H(x_i)$. The first approach selects points where the entropy exceeds a predefined global threshold, as demonstrated in~\autoref{fig:patching}. The second method flags locations where the entropy is significantly higher compared to the preceding value. This latter strategy can also be viewed as detecting instances where the entropy deviates from an otherwise roughly monotonically decreasing trend within the patch.

\begin{align*}
    \text{Global Constraint}\hspace{1cm}H(x_t)& > \theta_g\\
    \text{Approx. Monotonic Constraint}\hspace{1cm}H(x_t)& - H(x_{t-1}) > \theta_r
\end{align*}

Patch boundaries are determined through an efficient preprocessing procedure that takes place in the dataloader stage. This approach contrasts with the method employed by~\citet{nawrot-etal-2023-efficient}, which utilizes a classifier to detect patch boundaries based on entropy estimates. In our experimental analysis~(\S\ref{section:experiments}), we evaluate and contrast these strategies for differentiating between bytes of low and high entropy.

\subsection{The Byte-Pair Encoding (BPE) Tokenizer and Incremental Patching}
\label{section:incremental-patching}
\label{section:bpe-patch}

A majority of contemporary llms, such as our chosen baseline Llama 3, employ subword tokenizers based on bpe~\citep{gage1994new, sennrich-etal-2016-neural}. Throughout this work, we define ``tokens'' as sequences of bytes selected from a \textit{finite} and predetermined vocabulary established before model training, distinguishing them from ``patches,'' which denote adaptively formed segments that do not rely on a fixed vocabulary. One important distinction between tokens and patches is that, when utilizing tokens, the model lacks direct exposure to the original byte-level representations.

A significant advancement provided by {\textsc BLT} compared to tokenization-based approaches is its reconfiguration of the balance between vocabulary size and computational resources. Conventional llms must contend with the fact that expanding the vocabulary leads to an increase in the average token length, which decreases the total number of decoding steps required; however, this simultaneously enlarges the output dimensionality of the model's final projection layer. As a result, tokenization-based methods have limited ability to flexibly adjust both token granularity and inference computation. For instance, in moving from Llama 2 to Llama 3, the average token length grows from 3.7 to 4.4 bytes, but this improvement comes at the expense of quadrupling the embedding table size.

During generation, {\textsc BLT} must determine at each point in the byte sequence whether it is positioned at a patch boundary, as this choice dictates if the Latent Transformer will be activated for further computation. Importantly, this decision must be made in isolation from upcoming parts of the sequence, which have not yet been produced. Therefore, the patching mechanism cannot rely on unseen future bytes when segmenting the sequence. More precisely, a patching strategy $f_p$ adheres to the criterion of incremental patching if it holds that:
$$f_p(\pmb{x}_{<i}) = f_p(\pmb{x})_{<i}$$
bpe is an example of a patching method that does not possess the incremental property: given the same prefix, its tokenization can vary according to the suffix, thereby failing to meet the condition above.\footnote{Insertion of a unique delimiter to indicate patch demarcations can enforce incrementality for bpe; however, this also lengthens the byte sequence.}

\section{{\textsc BLT} Architecture}
\label{section:architecture}
The architecture of {\textsc BLT} includes a large-scale global autoregressive language model that operates over patch-level representations. Complementing this core component are two compact local models: one transforms byte sequences into patch representations, and another converts these patches back into byte sequences~(\autoref{fig:arch_high_level}).

\subsection{Latent Global Transformer Model}

\textit{The Latent Global Transformer} refers to an autoregressive transformer architecture denoted by $\mathcal{G}$, composed of $l_{\mathcal{G}}$ layers. This model transforms a series of latent input patch representations, $p_j$, into a corresponding series of output patch representations, $o_j$. In the notation used within this work, $j$ indexes patches, while $i$ indexes bytes. The global transformer leverages a block-causal attention mask~\citep{dubey2024llama}, ensuring that each token attends only to its own and preceding patches within the same document. Since this model accounts for the majority of computational resources required during both pre-training and inference, the frequency and context in which it is applied provides a mechanism to dynamically allocate computation, depending on the complexity of various segments of the input or output.

\subsection{Local Encoder}
\label{section:local}

\textit{The Local Encoder Model}, represented as $\mathcal{E}$, is designed as a compact, transformer-based architecture consisting of $l_{\mathcal{E}}$ layers, with $l_{\mathcal{E}}$ significantly smaller than $l_{\mathcal{G}}$. Its core function is to rapidly convert sequences of input bytes, $b_i$, into informative patch representations, $p_j$. Notably, this model deviates from standard transformer setups by integrating a cross-attention layer subsequent to each transformer layer; this module aggregates byte-level features to produce patch representations~(\autoref{fig:crossattn}).

Initially, input bytes $b_i$ are converted into vector representations $x_i$ through a learned embedding matrix of shape $\mathbb{R}^{256 \times h_{\mathcal{E}}}$. These vector embeddings may be further enriched with supplemental data using hash-embeddings~(\S\ref{sec:hash-emb}). The model then employs a stack of transformer and cross-attention layers in alternating fashion~(\S\ref{sec:enc-cross-attn}), which iteratively shape the byte-level inputs into the patch representations $p_i$ that are subsequently passed to the global transformer $\mathcal{G}$.

The attention mechanism within the transformer layers follows a \textit{local block causal} masking approach: each input byte may attend to a fixed window of $w_{\mathcal{E}}$ immediately preceding bytes. While this allows for information flow across patch boundaries, it does not permit attention across different documents. The subsequent sections elaborate on the structure of the embeddings and provide a closer examination of the cross-attention module.

\subsubsection{Encoder Hash n-gram Embeddings}
\label{sec:ngram-emb}
\label{sec:hash-emb}
An essential aspect of producing strong and informative representations at each position $i$ is to effectively capture context from prior bytes. {\textsc BLT} addresses this challenge by representing the current byte $b_i$ not only on its own, but also within the context of surrounding bytes as an n-gram. At every step $i$, we begin by generating byte-grams

\begin{align}
    g_{i,n}=\{b_{i-n + 1},\ldots, b_i\}
\end{align}

for each byte position $i$ and $n$ from three to eight.\footnote{
We omit byte-grams of size $n$ or more when $i<n$. }

Next, we present hash $n$-gram embeddings, where each byte $n$-gram is assigned to an index in a fixed-size embedding table $E_{n}^{hash}$ through the use of a hash function, for each $n$ in $\{3,4,5,6,7,8\}$~\citep{bai2010learning}. The generated $n$-gram embedding is subsequently summed with the corresponding byte's embedding, after which it undergoes normalization and is supplied as input to the local encoder module. The augmented embedding is computed

\begin{align}
e_i &= x_i + \sum_{n = 3,...,8}E_{n}^{hash}(\text{Hash}(g_{i,n})) \\
\text{where, Hash}(g_{i,n}) &= \text{RollPolyHash}(g_{i,n}) \% |E_{n}^{hash}|
\end{align}

We scale $e_i$ by dividing it by the total number of $n$-gram sizes incremented by one, employing RollPolyHash as described in Appendix~\ref{appendix:rollpolyhash}. Section~\ref{section:ablations} presents an ablation study examining how variations in $n$-gram size and the embedding table dimension for $n$-gram hash embeddings influence trends in flop-controlled scaling laws. Beyond hash-based $n$-gram embeddings, we have also explored $n$-gram embeddings determined by frequency, with methodological specifics included in Appendix~\ref{appendix:ngrams}.

\subsubsection{Encoder Multi-Headed Cross-Attention}
\label{sec:enc-cross-attn}
Our implementation is largely based on the input cross-attention component described in the Perceiver architecture~\citep{jaegle2021perceiver}. The primary distinction lies in our use of variable patch representations as the latent entities, rather than utilizing a static group of latent vectors as done originally~(\autoref{fig:crossattn}). Additionally, our latent representations are restricted to attending solely to the byte sequence forming their respective patch. Each patch $p_j$ is associated with a unique query vector, generated by first aggregating the byte-level representations within patch $p_j$ through a pooling operation, followed by application of a linear transformation, $\mathcal{E}_{C} \in \mathbb{R}^{h_{\mathcal{E}} \times (h_{\mathcal{E}} \times U_{\mathcal{E}})}$. Here, $U_{\mathcal{E}}$ denotes the number of attention heads within the encoder cross-attention mechanism. Let $f_{\text{bytes}}(p_j)$ be the ordered set of bytes constituting the patch $p_j$; the calculation proceeds as:

\begin{align}
P_{0,j} &= \mathcal{E}_{C}(f_\text{bytes}(p_j)), \quad f~\text{is a pooling function} \\
P_l &= P_{l-1} + W_o\left(\text{softmax}\left(\frac{QK^T}{\sqrt{d_k}}\right)V\right) \\
\text{where } Q_j &= W_q(P_{l-1,j}), K_i = W_k(h_{l-1,i}), V_i = W_v(h_{l-1,i}) \\
h_l &= \text{Encoder-Transformer-Layer}_l(h_{l-1})
\end{align}

Here, $P \in \mathbb{R}^{n_p \times h_{\mathcal{G}}}$ encodes $n_p$ patch representations intended for processing by the global model, initially formed by aggregating the byte embeddings $e_i$ associated with each patch $p_j$. The matrices $W_q$, $W_k$, $W_v$, and $W_o$ denote the linear transformations applied to queries, keys, values, and outputs, with the keys and values specifically derived from transforming the byte representations $h_i$ obtained from the preceding layer (or $e_i$ in the initial layer). A patch-specific masking approach is adopted, ensuring that each query $Q_j$ is restricted to attending exclusively to the keys and values representing the bytes within its own patch $j$. As multi-head attention is utilized over $Q, K,$ and $V$, and because patch representations generally possess greater dimensionality ($h_{\mathcal{G}}$) compared to $h_{\mathcal{E}}$, $P_l$ is structured as several heads each of dimension $h_{\mathcal{E}}$ for cross-attention operations; these heads are subsequently concatenated to yield vectors of size $h_{\mathcal{G}}$. A pre-LayerNorm is applied to queries, keys, and values, and no positional encodings are incorporated in this cross-attention mechanism. Additionally, a residual connection is introduced around the cross-attention layer.

\subsection{Local Decoder} 

The local decoder $\mathcal{D}$, akin to the local encoder, is a lightweight transformer architecture characterized by $l_{\mathcal{D}} << l_{\mathcal{G}}$ layers. Its purpose is to map the series of global patch representations, $o_j$, back into the raw byte sequence, $y_i$. In generating raw bytes sequentially, it leverages the history of already generated bytes, utilizing the hidden states generated by the local encoder for the corresponding byte sequence as input. The architecture consists of $l_{\mathcal{D}}$ layers alternating between cross-attention and transformer blocks. Within each layer, cross-attention precedes the transformer operation, effectively translating patch-level representations to byte-level ones, after which the transformer's local decoder layer processes this newly-formed byte sequence.

\subsubsection{Decoder Multi-headed Cross-Attention} 

Within the decoder’s cross-attention module, the query and key/value assignments are inverted; that is, byte representations now serve as the queries, while the patch representations assume the role of keys and values. At the outset, byte representations for cross-attention are set to the byte embeddings output from the final encoder layer, specifically $h_{l_{\mathcal{E}}}$. For each decoder layer $l$, the byte representation $d_{l,i}$ is updated using the following procedure:

\begin{align}
D_0 &= h_{l_{\mathcal{E}}} \\
B_l &= D_{l-1} + W_o\left(\text{softmax}\left(\frac{QK^T}{\sqrt{d_k}}\right)V\right), \\
  &\text{where } Q_i = W_q(d_{l-1,i}), K_i = W_k(\mathcal{D}_C(o_j)), V_i = W_v(\mathcal{D}_C(o_j)) \\
  D_l &= \text{Decoder-Transformer-layer}_l(B_l)
\end{align}

Here, as previously, $W_k, W_v$ denote the key and value projection matrices, responsible for applying a linear transformation followed by a split operation $\mathcal{D}_C$ to the final patch representations $o_j$ generated by the global model. The query projection matrix $W_q$ is applied to the byte-level representations $d_{l-1}$ obtained from the preceding decoder transformer layer (or $h_{l_{\mathcal{E}}}$ when processing the initial layer). $W_o$ serves as the output projection matrix. Consequently, $B$ has dimensions $\mathbb{R}^{h_{\mathcal{D}} \times n_b}$, with $n_b$ representing the number of bytes in the output. The succeeding decoder representations, $D_l$, are produced by running a decoder transformer layer on the output $B$ from the cross-attention module. Mirroring the approach in the local encoder cross-attention, we employ multi-head attention, incorporate pre-LayerNorms, exclude positional embeddings, and utilize a residual connection that encloses the cross-attention operation.

\section{Experimental Setup}
\label{section:experiments}
We construct a series of meticulously controlled experiments in order to rigorously compare {\textsc BLT} against models that rely on tokenization, explicitly ensuring that {\textsc BLT} does not benefit from potentially longer context windows.

\subsection{Pre-training Datasets}  
All model variants analyzed in this study are pre-trained on two distinct corpora: 1) The Llama 2 dataset~\citep{touvron2023llama}, which consists of 2 trillion tokens amassed from diverse publicly accessible resources, followed by extensive cleaning and filtering procedures to enhance dataset integrity; and 2) {\textsc BLT}-1T: An original dataset containing 1 trillion tokens, sourced from an array of public datasets, as well as incorporating a portion of the pre-training data made available by Datacomp-LM~\citep{li2024datacomp}. The Llama 2 dataset is leveraged in scaling law investigations, specifically to determine the optimal token quantity in accordance with the findings of~\cite{dubey2024llama} for guiding the architectural decisions in {\textsc BLT}, whereas {\textsc BLT}-1T serves as the basis for comprehensive pre-training runs that enable comparisons to Llama 3 across downstream benchmarks. Neither corpus contains data originating from Meta's own products or services. In addition, for evaluations involving tokenizer-centric baselines, we implement the Llama 3 tokenizer, characterized by a 128K token vocabulary, as this configuration yielded superior results compared to the Llama 2 tokenizer during our experiments.

\subsection{Entropy Model}  
For our experiments, the entropy model consists of a byte-level language model, which is trained using the identical training data distribution as the comprehensive {\textsc BLT} model. By default, unless specified otherwise, this model employs a transformer architecture comprising 100 million parameters, 14 layers, a hidden dimension size of 512, and a sliding window attention span covering 512 bytes. All other hyperparameter settings match those of the local and global transformer configurations utilized elsewhere in our approach. Various configurations were evaluated—including model scale, receptive field size, and network structure—as detailed in \autoref{section:ablations}. Notably, with sufficiently limited receptive fields, the resultant entropy model lends itself to representation as an optimized lookup table.

\subsection{Entropy Threshold and Equalizing Context Length} For approaches leveraging entropy-based patching, we determine a patching threshold that yields a targeted mean \textit{patch size} across the pretraining data mixture. Within {\textsc BLT}, the selection of \textit{patch size} is not constrained as in token-based systems, and this flexibility significantly affects the model's context window. To ensure comparability and prevent larger patches from obtaining disproportionate benefits, we enforce an equivalent expected total number of bytes per batch. Consequently, for models employing larger patch sizes, sequence length is curtailed accordingly. For models trained on Llama 2 data, the context window is set to 8k bytes, whereas for those trained on the {\textsc BLT}-1T dataset, the average context length is increased to 16k bytes—while the expected batch size remains at 16M bytes.

Although the batch size remains fixed on average, utilizing dynamic patching techniques leads to varying proportions of bytes per patch when assembling data batches. To enhance efficiency, our approach to {\textsc BLT} training compiles batches based on patch count, which eliminates the need for padding within the computationally costly latent transformer, thereby guaranteeing that each batch consists of an identical number of patches. During the training process, we pad or occasionally truncate the byte sequences to a maximum of 12k for the Llama 2 dataset and 24k for the {\textsc BLT}-1T dataset. This is done to mitigate memory usage surges that might arise from sequences containing exceptionally large patches.

\subsection{Entropy Model Context}
\label{section:entropy-context}
Our observations indicate that entropy patching tends to generate increasingly larger patches when applied to highly structured datasets, such as multiple choice problems (refer to the MMLU case in \autoref{fig:mmlu_patching}), where significant repetition occurs. The variation in patch sizes is attributable to reduced entropy within sections of repetitive content as detected by the entropy model context. Consequently, for the extensive execution of {\textsc BLT}-Entropy with a patch size of 4.5, we opted to refresh the entropy context using new line markers and employed an approximate monotonicity constraint, as this method is less impacted by "entropy drift" related to fluctuations in context length. Notably, this modification only influences the manner in which entropy values are determined; our approach for setting the entropy threshold remains unchanged.

\subsection{FLOPs Estimation} 
\label{section:flops}

Our calculation of transformer floating-point operations (flops) is primarily based on the methodology described in Chinchilla~\citep{hoffmann2022training}, which accounts for flops in the feed-forward components, qkvo projections within the self-attention module, as well as the attention computation and output projection steps. However, unlike Chinchilla, we assume that the input embedding layer operates through an efficient lookup mechanism rather than via dense matrix multiplication, rendering it a 0-flop step in our model. Consistent with established conventions, we assume that the backward propagation stage entails twice the flops of the forward propagation.

To determine the flops \textit{per byte} required by {\textsc BLT} models, we aggregate the flops from several components: the local encoder transformer, the global latent transformer, and the local decoder transformer, as well as the contributions from cross attention modules present in both the encoder and decoder. This aggregation is formally expressed as:

\begin{align}
    \text{FL}_{\text{{\textsc BLT}}} &= \text{Transf. FL}(h_{\mathcal{G}}, l_{\mathcal{G}}, m=n_{ctx}/n_p, V=0) / n_p\\
   &+\text{Transf. FL}(h_{\mathcal{E}}, l_{\mathcal{E}}, m=w_{\mathcal{E}}, V=0) \\
   &+ \text{Transf. FL}(h_{\mathcal{D}}, l_{\mathcal{D}}, m=w_{\mathcal{D}}, V=256) \\ 
    &+ \text{Cross Attn. FL}(h_{\mathcal{E}}, l_{\mathcal{E}}, m=n_p, r=n_p/k) \times k / n_p \\ 
   &+ \text{Cross Attn. FL}(h_{\mathcal{D}}, l_{\mathcal{D}}, m=k, r=k/n_p) 
\end{align}

Here, $n_{ctx}$ denotes the length of the sequence measured in bytes, $n_p$ represents the size of the patch, $r$ indicates the proportion of queries relative to key/values, and $k$ denotes the ratio of patch dimension to byte dimension; specifically, $k$ is the number of separate local model segments combined to assemble the overall model representation (for example, $k=2$ in \autoref{fig:crossattn}). $V$ specifies the size of the output projection's vocabulary, which is used exclusively in the local decoder. The attention mechanism adopts different context lengths, $m$, depending on whether the module is processing the sequence as bytes or as patches. We update the formulas for attention flops for each module accordingly. Detailed formulae for computing Transformer-FLOPs and Cross-Attention FLOPs can be found in Appendix~\ref{section:flops-appendix}.

\subsection{Bits-Per-Byte Estimation} Perplexity is meaningful solely within the framework of a specific tokenizer, as it quantifies uncertainty per token. To facilitate fair comparison between models operating on bytes and those utilizing tokens, we adopt Bits-Per-Byte (BPB)—a tokenizer-agnostic analogue to perplexity—as recommended in prior research~\citep{xue2022byt5, yu2023megabyte, wang2024mambabyte}. In detail:

\begin{align}
    \text{BPB}(x)&= \frac{\mathcal{L}_{CE}(\pmb{x})}{\ln(2)\cdot n_{\text{bytes}}}
\end{align}

where the uncertainty over the data $\pmb{x}$ as measured by the sum of the cross-entropy loss is normalized by the total number of bytes in $\pmb{x}$ and a constant. 

\subsection{Transformer Architecture Hyperparameters} 

In configuring all transformer blocks within {\textsc BLT}, which includes both local and global models, our approach aligns closely with the Llama~3 framework~\citep{dubey2024llama}. The feed-forward layers make use of the SwiGLU activation function~\citep{swiglu}, while rotary positional embeddings (RoPE)~\citep{RoPe} with parameter $\theta=500000$~\citep{xiong2024effective} are incorporated exclusively within the self-attention modules. Additionally, RMSNorm~\citep{RMSNorm} is employed for layer normalization throughout the architecture. All self-attention operations utilizing fixed-standard attention masks (e.g., \textit{block causal} or \textit{fixed-window block causal}) leverage Flash attention~\citep{dao2022flashattention}, with a window size set at 512 for fixed-width attention patterns. For cross-attention components—where dynamic, patch-specific masks are used—we employ Flex Attention\footnote{\url{https://pytorch.org/blog/flexattention}} to facilitate fused operations, resulting in a substantial acceleration of training times.

\subsection{{\textsc BLT}-Specific Hyperparameters}  
To evaluate the performance of {\textsc BLT} models, our experiments focus on two main aspects: analyzing scaling behaviors and assessing performance on downstream tasks, considering model sizes of 400M, 1B, 2B, 4B, and 8B parameters. The hyperparameters describing the architectures for these different scales are detailed in Appendix~Table~\ref{tab:arch}. 

Initialization of queries for the first cross-attention layer within the local encoder is performed through max-pooling. Across all {\textsc BLT} models, we utilize a single hash function generating $500,000$ hashes, with n-grams sizes that vary between 3 and 8. Every model is trained using a learning rate set at $4e-4$. Selection of this learning rate, for both token-based and {\textsc BLT} architectures, results from a hyperparameter sweep in the $1e-3$ to $1e-4$ range on 400M and 1B scale models, identifying $4e-4$ as optimal for both cases. For scaling experiments employing Llama-2 data, batch sizes conform to recommendations from~\cite{dubey2024llama} or are matched by equivalent memory usage in bytes. AdamW~\citep{adamw} serves as the optimizer, configured with $\beta_1 = 0.9$, $\beta_2 = 0.95$, and $\epsilon=10^{-8}$. We implement a linear warm-up phase for the first 2000 iterations, followed by cosine decay of the learning rate toward zero; weight decay is fixed at 0.1 and we enforce a maximum global gradient norm of 1.0 through gradient clipping.

\section{Scaling Trends}
\label{section:scaling}

We provide a comprehensive overview of the scaling dynamics of byte-level models, offering guidance for future scaling of {\textsc BLT} architectures. Our investigation into scaling addresses gaps in prior studies of byte-level models through several key contributions: (a) We analyze performance under compute-optimal training conditions, (b) We train 8B parameter models on substantial datasets—reaching up to 1T tokens or 4T bytes—and assess their effectiveness on a range of downstream benchmarks, and (c) We examine scaling patterns with respect to fixed inference costs. A subsequent section is dedicated to exploring the distinct benefits associated with modeling data as byte sequences.

\subsection{Parameter Matched Compute Optimal Scaling Trends}
\label{section:parameter-matched-scaling}
We conduct experiments on the Llama 2 dataset by training a range of \textit{compute-optimal} bpe and {\textsc BLT} models with parameter counts spanning from 1B to 8B. For each configuration, we chart the training FLOPs against language modeling accuracy over a representative slice of the training data mixture. The bpe models leverage the theoretically ideal model-to-data ratio, as indicated by Llama~3~\citep{dubey2024llama}, to ensure optimal utilization of computational resources as established in prior work~\citep{hoffmann2022training}. This approach yields a strong baseline, maximizing performance relative to computational budget. In parallel, every bpe model has an identically sized {\textsc BLT} counterpart, each trained on the same dataset with a Latent Transformer mirroring the architectural choices of the paired bpe Transformer.

As shown in~\autoref{fig:scaling-compute-opt} (right), the {\textsc BLT} models consistently perform at least as well as, and often surpass, their bpe-based counterparts, a pattern that remains consistent as both model size and computational resources are scaled up. To our knowledge, this marks the first instance of a byte-level Transformer architecture, {\textsc BLT}, demonstrating scaling behavior comparable to that of BPE-based models under compute-optimal conditions. This observation supports our hypothesis that the parameter-to-training-compute ratio identified as optimal for bpe models is also applicable to {\textsc BLT}, or deviates from it only minimally.

Enhancements at the architectural level, in tandem with the adoption of dynamic patching, are essential for achieving favorable bpe scaling behavior. As depicted in ~\autoref{fig:scaling-compute-opt} (left), we conduct a comparison between space-patching models and Llama 3. Here, we approximate SpaceByte \citep{slagle2024spacebyte} by utilizing {\textsc BLT} space-patching while omitting n-gram embeddings and cross-attention mechanisms. While SpaceByte shows gains relative to Megabyte, there is still a considerable gap to Llama 3. In \autoref{fig:scaling-compute-opt} (right), we demonstrate the contributions from both architectural refinements and the application of dynamic patching. At larger model scales, {\textsc BLT} variants achieve performance comparable to leading tokenizer-based models, including Llama 3.

Additionally, we note that performance in tokenizer-based models is influenced by the specific tokenizer employed; in particular, models trained with the Llama-3 tokenizer consistently outperform those trained with the Llama-2 tokenizer when evaluated on identical data.

In summary, the {\textsc BLT} model demonstrates behavior that falls between that of Llama 2 and 3, especially when employing much larger patch sizes. While the bpe tokenizers for Llama 2 and 3 have average token lengths of 3.7 and 4.4 bytes, respectively, {\textsc BLT} can mirror similar scaling dynamics using average patch sizes of 6 or even 8 bytes. Since inference flops are inversely related to average patch size, choosing an 8-byte patch size offers nearly a 50\% reduction in inference computational cost. Additionally, models configured with larger patch sizes appear to improve their performance as both the model and data scale up. Notably, {\textsc BLT} with an 8-byte patch size initially performs considerably worse than bpe Llama 2 at the 1B parameter scale but surpasses bpe performance at the 7B scale. These outcomes indicate that larger patch sizes may yield superior results as models and datasets increase further in scale, and it is conceivable that even greater patch sizes could be advantageous as training resources expand.

\subsection{Beyond Compute Optimal Task Evaluations}
\label{section:evals}
To further analyze scaling behaviors, we train an 8B {\textsc BLT} model at a point exceeding the compute optimal regime, utilizing the {\textsc BLT}-1T dataset—characterized by its increased size and improved data quality—and evaluate its capabilities across a variety of established classification and generation benchmarks. The evaluation suite encompasses common sense reasoning, world knowledge, and code generation tasks, detailed as follows:
\paragraph{Classification tasks} consist of ARC-Easy (0-shot)~\citep{Eval_ARC}, Arc-Challenge (0-shot)~\citep{Eval_ARC}, HellaSwag (0-shot)~\citep{Eval_hellaswag}, PIQA (0-shot)~\citep{Eval_piqa}, and MMLU (5-shot)~\citep{hendrycks2020measuring}. For these, we utilize a prompt-scoring framework that computes the model's predicted probability for each answer option and calculate the average accuracy across tasks.
\paragraph{Coding related generation tasks:}  We evaluate Python code synthesis using the pass@1 metric on MBPP (3-shot)~\citep{Eval_MBPP} and HumanEval (0-shot)~\citep{Eval_HumanEval}, which reflects the proficiency of LLMs in generating correct code in a single attempt.

In \autoref{tab:evals}, we present a comparative analysis of three models, all trained on the {\textsc BLT}-1T dataset: a model utilizing the bpe Llama 3 tokenizer,\footnote{We opt for the Llama 3 tokenizer and its 128k vocabulary, as it yields superior results relative to the 32k vocabulary of Llama 2.} and two distinct configurations of the {\textsc BLT} architecture. One variant leverages a space-based patching method ({\textsc BLT}-Space), while the other employs a patching technique guided by entropy values ({\textsc BLT}-Entropy). The latter incorporates an approximate monotonicity constraint and resets its context whenever a new line occurs, in accordance with the procedure described in~\autoref{section:entropy-context}. Training for all models is conducted with an identical computational (flop) budget. Notably, for {\textsc BLT}-Entropy, we adjust the entropy threshold at inference time from 0.6 to 0.1, an alteration that enhances performance on evaluated tasks though it comes at the expense of increased inference steps.

The {\textsc BLT}-Entropy model achieves superior results compared to the Llama 3 model on four of the seven evaluated tasks, despite both being trained on an equivalent volume of data. This enhancement can likely be attributed to two key factors: (1) more efficient utilization of computational resources during training through dynamic patching, and (2) the explicit handling of byte-level data rather than token-level representations.

In contrast, {\textsc BLT}-Space generally lags behind the Llama 3 tokenizer across nearly all tasks except one. Nevertheless, it realizes considerable computational savings during inference owing to its increased mean patch size of 6 bytes. For comparison, models utilizing bpe or entropy-based patching display an average patch size close to 4.5 bytes when assessed on the training data distribution. Given identical training resource constraints, the model with the larger patch size is able to process 30\% more data than its counterparts, potentially shifting {\textsc BLT} even further from the compute-optimal regime.

\subsection{Patches Scale Better Than Tokens} Utilizing {\textsc BLT} models enables a concurrent scaling up of both model parameters and patch dimensions, all while preserving the computational cost for training and inference and maintaining a fixed training dataset. This flexibility in expanding patch size distinguishes patch-based models from conventional token-based approaches, which are constrained by the limitations of a static vocabulary, as elaborated in Section~\ref{section:bpe-patch}. Employing larger patches reduces computational requirements, thereby allowing resources to be redirected toward expanding the global latent transformer, since it operates at a reduced frequency.

A fixed inference scaling study is performed to evaluate whether increasing model size, combined with processing fewer, larger patches, results in superior performance compared to deploying smaller models that operate on more patches. Using initial models of 400m and 3.6B parameters with the Llama 2 tokenizer, we identify models with comparable computational cost using both the Llama 3 tokenizer and {\textsc BLT}-Entropy models, the latter employing average patch sizes of 6 and 8 bytes within the training datamix (refer to~\autoref{tab:fixed-inf-params} for additional model specifications). For models with an 8-byte patch size, 3 encoder layers are utilized rather than just 1. Each architecture is trained under varying training flop budgets.

As illustrated in \autoref{fig:fixed_inference_scaling}, {\textsc BLT} models display superior scaling behavior compared to tokenization-based architectures within both inference flop categories. While BPE models initially exhibit stronger performance under limited training resources, their advantage diminishes as training budgets increase, with {\textsc BLT} models overtaking them shortly past the compute-optimal threshold. In many scenarios, allocating additional resources for the pretraining phase may be advantageous in order to obtain models that deliver enhanced performance given a fixed inference computation. The 8B parameter models, such as Llama 3.1, exemplify this principle; these models have been pretrained on data volumes nearly two orders of magnitude greater than the compute-optimal amount corresponding to their size.

The point at which {\textsc BLT} surpasses token-based models now occurs at a compute scale somewhat closer to the compute-optimal threshold when transitioning to models with a larger flop class, decreasing from approximately 3x to 2.5x the compute-optimal allocation. Likewise, the patch size 8 variant exhibits a more rapid scaling progression within the larger flop class, overtaking competing architectures earlier. As elaborated in Section~\ref{section:parameter-matched-scaling}, it is evident that at increased model sizes, larger patch sizes yield performance more comparable to BPE models. This observation can be partially ascribed to a reduction in the proportion of total flops expended on the byte-level Encoder and Decoder, which exhibit slower scaling relative to the Latent Transformer component. When scaling the overall parameter count by a factor of 20—from 400M up to 8B—there is only an approximate twofold increase in {\textsc BLT}'s local model parameters. This distinction is crucial because enlarging the patch size influences only the flops for the patch Latent Transformer, leaving those for the byte-level modules unaffected. Consequently, the ratio of {\textsc BLT}-Entropy ps=8 parameters to those of the Llama 2 model rose from 1.6x to 1.7x as model scale increased.

In conclusion, our investigation into scaling with respect to patch length shows that the {\textsc BLT} patch-based architecture benefits from concurrent increases in both patch size and model capacity, resulting in superior scaling behavior. This positive scaling pattern appears to not only endure but even strengthen at larger model sizes.

\section{Byte Modeling Improves Robustness} We further evaluate the robustness of {\textsc BLT} against token-based models that do not access explicit byte-level data, and describe a method for converting pretrained token-based models to operate at the byte level.
\label{sec:robustness}

\subsection{Character-Level Tasks}

One of the initial driving factors behind the development of byte-level models was their capacity to handle noise introduced at the byte level within inputs, as well as their intrinsic understanding of token composition—areas where traditional models relying on tokenizers often face challenges. To assess these specific properties, we conducted supplementary experiments using benchmarks designed to test models' resilience against input noise and their sensitivity to character-level information, spanning English and multiple languages, and also encompassing digits and phonemes. The outcomes of these assessments are summarized in Table~\ref{tab:char_tasks}.

\paragraph{Noisy Data} To assess and contrast the resilience of tokenizer-based models and {\textsc BLT}, we generate altered versions of the benchmark classification tasks outlined in Section~\ref{section:evals} by introducing character-level noise. We utilize five specific character-based noising methods to modify the input text: (a) \textit{AntSpeak}, where the entire text is rendered in uppercase with spaces inserted between each character; (b) \textit{Drop}, which involves randomly eliminating 10\% of the characters; (c) \textit{RandomCase}, in which half of the characters are randomly made uppercase and the other half lowercase; (d) \textit{Repeat}, where 20\% of the characters are duplicated up to four times; and (e) \textit{UpperCase}, which simply converts all text to uppercase letters. For evaluation, each noising technique is individually applied either to the prompt, the completion, or both, treating each scenario as a separate task, and we compute the mean performance. Table \ref{tab:char_tasks} presents the outcomes for noised HellaSwag~\citep{Eval_hellaswag}, demonstrating that {\textsc BLT} consistently exhibits greater robustness relative to tokenizer-based counterparts, achieving an average margin of 8 points over an equivalently trained model, and surpassing the Llama 3.1 model, despite its substantially larger training corpus.

\paragraph{Phonology - Grapheme-to-Phoneme (G2P)} We evaluate how effectively {\textsc BLT} converts sequences of graphemes, which are the written symbols of a word, into phonemes that represent its pronunciation. Table \ref{tab:char_tasks} displays outcomes for the G2P task under a 5-shot regime utilizing Phonology Bench~\citep{suvarna2024phonologybench}. Our findings indicate that {\textsc BLT} achieves superior performance compared to the baseline Llama 3 1T model using tokenizer-based approaches.

\paragraph{CUTE}
To evaluate the extent of character-level comprehension, we measure {\textsc BLT}'s performance on the CUTE benchmark~\citep{edman2024cute}, which includes a variety of tasks grouped into three main categories: compositional understanding, orthographic similarity, and sequence manipulation capabilities. This suite of tasks presents substantial difficulty for most models reliant on tokenization methods; although these models recognize the orthographic structure of their tokens, they often fail to leverage this knowledge effectively for text manipulation. As reported in Table~\ref{tab:char_tasks}, the {\textsc BLT}-Entropy model surpasses both BPE Llama 3 variants on this benchmark by over 25 points. Our model exhibits particularly strong aptitude in tasks requiring fine-grained manipulation at the character level, achieving a 99.9\% score on both spelling tasks. These significant performance gains, attained even though {\textsc BLT} is trained on sixteen times less data than Llama 3.1, suggest that mastering character-level detail remains a persistent challenge for models built on BPE tokenization. Figure \ref{fig:cute_examples} showcases cases where the Llama 3 tokenizer-based model encounters difficulties, whereas {\textsc BLT} achieves accurate results. The only areas where BPE models hold a relative advantage are in tasks involving word deletion and insertion. While such manipulations might present inherent challenges for a byte-level architecture, the performance gap remains moderate, and composing word-level understanding from characters might, in fact, be more tractable than the reverse. Throughout, we employ consistent evaluation methodologies and use original prompts sourced from Huggingface. It should be noted that BPE-based models could potentially realize gains from further prompt optimization.

\paragraph{Low Resource Machine Translation} We assess the performance of {\textsc BLT} for both translation into and out of twenty-one low-resource languages across six widely recognized language families. These languages, using a variety of scripts, are drawn from the FLORES-101 benchmark~\citep{goyal2022flores}. SentencePiece BLEU scores are summarized in Table~\ref{tab:flores}. The evaluation indicates that {\textsc BLT} exceeds the performance of a baseline system utilizing the Llama 3 tokenizer, offering an overall improvement of 2 BLEU points for translations into English, and a 0.5-point improvement for translations out of English. For established language pairs, {\textsc BLT} matches or slightly surpasses Llama 3. Importantly, for numerous language pairs involving underrepresented language families, {\textsc BLT} shows a marked advantage over Llama 3, highlighting the benefits of byte-based modeling for capturing long-tail byte sequence generalization in low-resource scenarios.

\subsection{Training {\textsc BLT} with Llama 3 Initialization}
\label{sec:distillation}
We investigate a method in which {\textsc BLT} models are able to benefit from previously pre-trained, tokenizer-based models to enhance training efficiency and speed. This is accomplished by initializing the global transformer parameters in {\textsc BLT} with those taken from a pre-trained Llama 3.1 model. After this initialization step, we further adjust the global transformer weights by applying a learning rate one-tenth the size of that used for the local encoder and decoder, training for the same number of steps recommended for Llama 3. The results, shown in Table~\ref{tab:distill}, allow for direct comparison with a baseline {\textsc BLT} approach. The findings demonstrate that {\textsc BLT} initialized from Llama 3.1 offers substantial gains over both baseline Llama 3 and baseline {\textsc BLT}, with all models trained under equivalent computational budgets. Furthermore, even in comparison with the {\textsc BLT}-Entropy variant—refer to Table~\ref{tab:evals}—which was trained on a much larger dataset (1T tokens), the {\textsc BLT} initialized from Llama 3.1 attains higher MMLU performance. These results indicate that this strategy offers a promising means for notably lowering training computation requirements.

This approach effectively transforms models that rely on tokenizers into models that do not, thereby adapting a pre-trained LLaMA 3.1 into a {\textsc BLT} architecture. For thorough evaluation, we report results for the original LLaMA 3.1, trained on 15T tokens, alongside the {\textsc BLT} variant initialized from LLaMA 3, as shown in Table \ref{tab:distill}. While our adapted model shows only a minor reduction in performance on MMLU and HumanEval, performance deteriorates more noticeably on other benchmarks. These findings indicate that additional research is required to more fully utilize the pre-trained model’s capabilities, especially by refining data mixture strategies and tuning hyperparameters.

\section{Ablations and Discussion}
\label{section:ablations}
This section presents ablation experiments that support our design decisions regarding the architecture of {\textsc BLT}, as well as details of the patching strategy and selection of hyper-parameters for the {\textsc BLT} model comprising 8B parameters, which was trained on the {\textsc BLT}-1T dataset.

\paragraph{Entropy Model Hyper-parameters} To investigate how the scale of the entropy model and the context window size influence performance, we train several byte-level entropy transformer models, varying their capacities from 1 million to 100 million parameters and utilizing context windows ranging from 64 to 512 tokens. We display the relationship between bits-per-byte (bpb) and training computational cost as scaling law curves, generated from our $400m$ and $1b$ parameter {\textsc BLT} models trained on the Llama-2 dataset, as shown in \autoref{fig:entropy_ablation}. Our observations indicate that increasing both the model size and the context window generally improves scaling performance, though the benefits plateau beyond 50 million parameters.

\paragraph{Types of Patching}

We conduct an ablation study on the four distinct patching strategies described in Section~\ref{section:patching}, namely: 1) Strided Patching utilizing strides of 4 and 6, 2) Whitespace-based Patching, 3) BPE Tokenizer patching employing the Llama 3 tokenizer, and 4) Entropy-driven patching via a compact byte-level language model.

Although dynamic patching shortens the effective sequence length, we standardize sequence length across all patching strategies to ensure a consistent context window. During both training and inference, each model is exposed to the same expected number of bytes per sequence. This approach eliminates potential confounds associated with varying context sizes. Figure~\ref{fig:scaling-compute-opt} presents the findings from these controlled experiments. The non-static patching strategies all show improvements over static patching, and notably, space patching performs nearly as well as dynamic entropy based patching.

\autoref{tab:evals_ablation} displays the results of benchmark evaluations for {\textsc BLT} models, offering a comparison among models that use standard tokenizers, space patching, and entropy-based patching. All models were trained on the Llama 2 dataset for an empirically selected number of optimization steps~\citep{dubey2024llama}. Despite the relative simplicity of space patching, which eliminates the need to compute entropy in real time during training, our results indicate that the benefits of entropy-based patching identified in our analysis of scaling behavior (see Section~\ref{section:scaling}) are preserved and evident in a range of downstream benchmark tasks.\footnote{Results for space patching are from previous experiments lacking cross-attention, but the same trends are apparent when cross-attention is included.}

\paragraph{Cross-Attention} 
In~\autoref{tab:crossattn_ablations_1b}, we examine the impact of introducing cross-attention at different locations within the encoder and decoder of {\textsc BLT} by conducting ablation studies. Specifically, for encoder cross-attention, we evaluate three approaches for initializing the queries: 1) employing an identical learned embedding for all global states, 2) using a hash embedding derived from the bytes present in the patch, and 3) applying pooling over the encoder’s hidden representations for the bytes within each patch at the relevant encoder layer.

Our results indicate that incorporating cross-attention into the \textit{decoder} yields the greatest benefit. When applied within the encoder, cross-attention offers a modest advantage, though this improvement is observed solely when queries are initialized through pooling. Moreover, the utility of cross-attention is especially pronounced for the Common-Crawl dataset and becomes more significant as patch sizes increase.

\paragraph{n-gram Hash Embeddings} 

We evaluate the impact of varying n-gram hash embedding vocabulary sizes—specifically, 0, 100K, 200K, and 400K—as reported in Table~\ref{tab:ngram_ablations_1b}. Our findings indicate that the inclusion of hash embeddings results in performance improvements across all evaluated domains, with especially pronounced effects for Wikipedia and Github, where the observed benefit is a 0.04 bpb improvement (relative to a 0.01 bpb improvement at 8B after 15k steps). At the 8B parameter scale, reducing the hash vocabulary from 500K to 300K resulted in only a marginal performance change of approximately 0.001 bpb at 15k steps. These results suggest that hash embeddings are crucial in enabling {\textsc BLT} models to achieve performance parity with those based on tokenizers; nonetheless, beyond 300K hashes, the advantage diminishes. Moreover, the evidence suggests that the improvements offered by hashes and cross-attention are mainly orthogonal, yielding benefits on distinct datasets.

\paragraph{Local Model Hyperparamaters} Table \ref{tab:local_layers} presents an ablation of different configurations for the number of layers in both the local encoder and decoder. Utilizing hash n-gram embeddings, {\textsc BLT} achieves strong performance with a very minimal encoder—specifically, using only a single layer—while benefitting from a more substantial, deeper decoder.

\section{Related Work}
\label{section:related_work}

\textbf{Character-Level RNNs:} Since the advent of neural networks, Character Language Modeling has garnered significant attention~\citep{sutskever2011generating,mikolov2012subword,graves2013generating}, largely because these models naturally handle out-of-vocabulary words without relying on explicit back-off strategies. In their work, \cite{kim2016character} develop a system that exclusively processes input at the character level via convolutional and highway layers, subsequently passing representations to LSTM-based RNNs; this approach achieves parity with contemporary RNN-based language models for English, and surpasses them for languages with complex morphology, showcasing a further strength of character-level LLMs. \cite{kenter2018byte} extend these ideas to machine comprehension, employing byte-level LSTM models which exceed the performance of word-based architectures particularly in morphologically rich languages such as Turkish and Russian. In parallel, \cite{zhang2015character} employ character-based convolutional frameworks for classification tasks and demonstrate that for some problems, these outperform word-level models. \citet{chung2019hierarchical} propose hierarchical LSTM models enhanced with boundary detectors at each layer to uncover textual hierarchies, further raising the ceiling on character-level language modeling performance. Unlike attention-centric architectures, ByteNet~\cite{kalchbrenner2016neural} applies convolutional operations to character sequences for tasks such as machine translation.

\textbf{Character-Level Transformers:} The advent of transformer architectures employing attention mechanisms~\citep{vaswani2017attention}, in conjunction with subword tokenization strategies~\citep{sennrich-etal-2016-neural}, has led to substantial advances in the accuracy of neural models for language modeling and standard benchmarks. Nevertheless, adopting word and sub-word representations instills a particular inductive bias, shaping the abstract level at which these models function. In an effort to harness both the advantages of transformers and the initial positive outcomes seen in character-level language modeling, \citet{al2019character} leveraged extremely deep transformer networks. Assisted by auxiliary loss functions, they succeeded in training transformer-based architectures that surpassed prior LSTM-based character-level language models. Despite these achievements, a noticeable performance disparity persisted when compared to word-level LLMs. Similarly, GPT-2~\citep{radfordgpt2} reported that, for large-scale datasets such as the 1 Billion Word Benchmark, byte-level language models still fell short of the results obtained by word-level counterparts.

Previous work by \cite{Choe2019BridgingTG} established that transformer-based LLMs operating at the byte level can surpass subword-based counterparts with similar parameter counts; however, these byte-level models require substantially greater computational resources and extended training times. In parallel, \cite{el2020characterbert} introduced CharFormer, a BERT variant that constructs word-level representations via convolutional operations applied to character embeddings. Although this method yielded notable improvements within the medical domain, it similarly incurred higher computational costs. Expanding on these lines, \cite{clark2022canine} proposed CANINE, a 150M-parameter encoder-only framework functioning directly on sequences of characters. CANINE incorporates a deep transformer stack, analogous to the global model architecture discussed here, while utilizing a hybrid design with a local transformer and strided convolutional layers to efficiently downsample the character inputs. This setup outperformed mBERT, a token-based encoder-only model, across multilingual downstream tasks. ByT5~\citep{xue2022byt5} investigated byte-level encoder-decoder configurations that forgo patching operations entirely. While this design promoted improved resilience to noisy input and achieved performance on par with traditional tokenizer-based models with only one-quarter of the data, omitting patching led to significant computational demands, as attention had to be computed for each byte. Handling bytes instead of subword units substantially lengthens input sequences, complicating the scalability of byte-level models. Most recently, leveraging the Mamba architecture~\citep{gu2023mamba}, characterized by its ability to maintain a constant-size memory state across long contexts, \cite{wang2024mambabyte} trained a byte-level Mamba model—again without patching—which surpassed byte-level transformers in flop-controlled experiments at the 350M parameter scale, achieving superior bits-per-byte performance over multiple datasets.

\textbf{Patching-based approaches:} Employing patching strategies can substantially reduce the high computational overhead typically associated with byte-level LLMs, potentially preserving their performance. Numerous studies have reported promising outcomes with these methods at modest model scales and data sizes. For instance, \cite{nawrot-etal-2022-hierarchical} explore the use of static patching through downsampling and upsampling operations, introducing the hourglass transformer that surpasses other byte-level approaches at the 150M parameter level. Subsequently, \cite{nawrot-etal-2023-efficient} enhance these results by applying dynamic patching techniques, which include an end-to-end trained boundary-predictor, a supervised boundary predictor guided by specific tokenizers, and an entropy-driven patching model comparable to {\textsc BLT}. They demonstrate that this methodology can exceed the performance of standard transformer architectures in language modeling tasks at the 40M parameter range, trained on roughly 400M tokens. Meanwhile, \cite{lester2024training} examine model training on sequences that have been compressed via arithmetic coding, achieving levels of compression that surpass those possible with BPE. Leveraging an equal-info windows method, they succeed in outperforming byte-level baselines in language modeling, though their results still lag behind subword-based models.

Our research is primarily motivated by MegaByte~\citep{yu2023megabyte}, a causal decoder-only large language model that leverages a fixed static patching approach, concatenating byte representations into patches, and employing a local model on the decoder to reconstruct bytes from these patches. The original MegaByte work shows that at the 1B parameter level, performance on a ~400B byte dataset is comparable to that of models relying on standard tokenizers. Throughout our experiments, we include extensive comparisons to MegaByte and observe that its static patching methodology falls short of the most recent compute-efficient tokenizer-based models in a flop-normalized context. We further illustrate how {\textsc BLT} overcomes this performance gap. In parallel, \cite{slagle2024spacebyte} report similar findings regarding MegaByte's limitations and propose augmenting static patching with alternative segmentations, such as using whitespaces and similar delimiters, in addition to incorporating a local encoder model. Their results suggest marginal gains over tokenized transformer models in a compute-matched regime on certain domains like Github and arXiv, specifically at the 1B parameter level. We further evaluate this model in our study and demonstrate that to reach, or surpass, the scaling capabilities of token-based systems like Llama 3, substantial additional architectural enhancements will be necessary for byte-level models.

\section{Limitations and Future Work}

For the purposes of selecting model architectures in this study, we adopt the optimal training steps established for Llama~3~\citep{dubey2024llama}. It should be noted, however, that these scaling laws are based on BPE-level transformer models and might not yield optimal relationships between data volume and parameter count when applied to {\textsc BLT}. Determining scaling laws specific to {\textsc BLT} remains an open question for future research and may reveal even more advantageous scaling behavior for this model class. Moreover, the majority of the present experiments have been limited to configurations with up to 1B parameters. As we scale models to 8B parameters and larger, it is conceivable that the best architectural choices will differ, potentially enabling greater performance at these increased scales.

Current transformer codebases are primarily optimized for high efficiency when working with tokenizer-based architectures. Although our experiments ensure theoretical flop equivalence and we employ various efficient solutions (like FlexAttention) to accommodate layers that differ from standard transformer designs, our implementations might still fall short of matching tokenizer-based models in wall-clock performance and could require additional optimization efforts.

Although {\textsc BLT} employs a separately pre-trained entropy model for its patching process, a promising avenue for future research would be to develop an end-to-end trained patching mechanism. As discussed in Section \ref{sec:distillation}, preliminary experiments demonstrate that it is feasible to apply ``byte-ifying'' techniques to tokenizer-based models such as Llama 3—which have been trained on datasets exceeding 10T tokens—by initializing and fixing the global transformer parameters using their original weights. Continued exploration along this trajectory could lead to approaches that not only preserve the advantages conferred by bytefying but also achieve superior performance relative to these tokenizer-based models, all without the necessity of fully retraining them.

\section{Conclusion}
\label{section:conclusion}

This work introduces the Byte Latent Transformer (\textbf{BLT}), an architectural innovation that moves beyond the traditional reliance on predefined token vocabularies in large language models. By employing a flexible, trainable approach to aggregating bytes into variable-length patches, {\textsc BLT} adapts resource allocation in response to data complexity, thereby enhancing both computational efficiency and model resilience. Through a comprehensive scaling analysis, we show that {\textsc BLT} models achieve parity with tokenization-based systems such as Llama 3 at scales up to 8B parameters and 4T bytes, and can reduce inference flops by up to 50\% with only slight compromises in standard evaluation metrics. Additionally, {\textsc BLT} introduces a novel axis for model scaling—enabling concurrent growth in model capacity and patch size without exceeding a fixed computational budget during inference. This approach proves beneficial for compute-constrained scenarios typical in applied machine learning contexts. By operating directly on byte sequences, {\textsc BLT} not only bolsters robustness to data noise and rare input cases, but also offers improved processing of sub-word information. Collectively, our results highlight {\textsc BLT} as a compelling alternative to conventional tokenization methods, promising greater scalability, robustness, and computational adaptability for large-scale language modeling tasks.

\end{document}